% 本文件由 tools/build_book.py 从 dist/reports/assertions.json 生成，请勿手改。
\begin{longtable}{@{}p{0.06\textwidth}p{0.32\textwidth}p{0.22\textwidth}p{0.16\textwidth}@{}}
\caption{章节 3.4 数值断言（生成时间 2026-10-19T16:08:50-0400）}\\
\toprule
类别 & 断言 & 理论值 & 实测值 \\
\midrule
\endfirsthead
\toprule
类别 & 断言 & 理论值 & 实测值 \\
\midrule
\endhead
\bottomrule
\endlastfoot

A & 帧契约：角点 $\times$ 时间端点的合法帧数 & 所有参数角点与时间端点都必须建出有限、\allowbreak{}合规的帧 & 4 个 viz，共 60 帧、\allowbreak{}711 个数组段，全部有限且 meta 全是有限标量 \checkmark \\
B & 帧内灵敏度 vs 直接重算(相对) & $\le$ 1.00000e-12 & 0 \checkmark \\
B & 灵敏度 vs 闭式 sinh(kL)/k(相对) & $\le$ 1.00000e-06 & $2.013\times 10^{-9}$ \checkmark \\
B & 帧内参考斜率 vs 直接重算(相对) & $\le$ 1.00000e-12 & 0 \checkmark \\
A & viz1 的数字对拍(decay, k=10, L=1) & S=\allowbreak{}sinh(\allowbreak{}kL)/\allowbreak{}k=\allowbreak{}1101.32；s\_ref=\allowbreak{}-k*coth(\allowbreak{}kL) & 帧内 S=\allowbreak{}1.101323e+03，直接重算 1.101323e+03(\allowbreak{}相对 0.0e+00)；闭式 1.101323e+03(\allowbreak{}相对 2.0e-09)；s\_ref=\allowbreak{}-10.000000 \checkmark \\
B & 帧内地板 vs 直接重算(相对) & $\le$ 1.00000e-12 & 0 \checkmark \\
B & 帧内半对数斜率/R\textasciicircum{}2 vs 直接重算 & $\le$ 1.00000e-09 & 0 \checkmark \\
A & viz2 的数字对拍(smooth) & 半对数斜率由未触地板的粗端窗口拟合；地板=\allowbreak{}曲线最小误差 & 斜率 -2.0894(\allowbreak{}重算 -2.0894)，R\textasciicircum{}2=\allowbreak{}0.99491，地板 7.550e-15(\allowbreak{}重算 7.550e-15)，窗口 N=\allowbreak{}4..16(\allowbreak{}head=\allowbreak{}7) \checkmark \\
B & 聚集比 vs 直接重算 & $\le$ 1.00000e-12 & 0 \checkmark \\
B & 密度比 vs 直接重算 & $\le$ 1.00000e-12 & 0 \checkmark \\
B & 密度比 vs 解析 1/sqrt(0.19) & $\le$ 1.00000e-12 & 0 \checkmark \\
A & viz3 的数字对拍(abs3, N=24) & 聚集比 0.6366N；密度比 1/\allowbreak{}sqrt(\allowbreak{}0.19)=\allowbreak{}2.294；配置解比插值差一阶 & 聚集比 15.2571(\allowbreak{}重算 15.2571)；密度比 2.294157(\allowbreak{}解析 2.294157)；插值峰值 7.76e-05，配置解误差 8.39e-03 \checkmark \\
B & ||D||\_\allowbreak{}inf vs N\textasciicircum{}2(相对) & $\le$ 1.00000e-12 & $6.661\times 10^{-16}$ \checkmark \\
B & 行和残差 vs 直接重算 & $\le$ 1.00000e-18 & 0 \checkmark \\
B & kappa+ vs 直接重算 & $\le$ 1.00000e-12 & 0 \checkmark \\
B & D\_\allowbreak{}00 vs (2N\textasciicircum{}2+1)/6(相对) & $\le$ 1.00000e-13 & $1.163\times 10^{-15}$ \checkmark \\
A & viz4 的数字对拍(N=16, D) & ||D||\_inf=\allowbreak{}N\textasciicircum{}2；D$\cdot$1=\allowbreak{}0；kappa+ \textasciitilde{} N\textasciicircum{}2；朴素 kappa 无意义 & ||D||=\allowbreak{}256.000000(\allowbreak{}N\textasciicircum{}2=\allowbreak{}256)，行和残差 4.44e-15，朴素 kappa=\allowbreak{}1.36e+18，kappa+=\allowbreak{}129.956 \checkmark \\
B & 灵敏度 vs sinh(kL)/k(相对，kL$\le$20) & $\le$ 1.00000e-08 & $1.653\times 10^{-9}$ \checkmark \\
B & kL=40 时灵敏度自身的相对偏差 & $\le$ 1.00000e-06 & $5.245\times 10^{-8}$ \checkmark \\
A & 变分方程灵敏度 vs 闭式 sinh(kL)/k & S =\allowbreak{} sinh(\allowbreak{}kL)/\allowbreak{}k & kL$\le$20 最大相对偏差 1.653e-09；kL=\allowbreak{}40 时 S=\allowbreak{}5.884631e+15(\allowbreak{}闭式 5.884632e+15，相对 5.2e-08) \checkmark \\
B & 线性问题 Newton 收敛后的斜率(与另一次打靶互差，相对) & $\le$ 1.00000e-08 & $5.764\times 10^{-16}$ \checkmark \\
A & 线性问题 Newton 迭代次数 & 1 步(\allowbreak{}仿射 F)+ 1 次确认求值 & 最多 2 次求值；smooth: s=\allowbreak{}-1.155727e+00，2 次求值，err=\allowbreak{}1.07e-10；neumann: s=\allowbreak{}1.570796e+00，2 次求值，err=\allowbreak{}2.47e-13；layer: s=\allowbreak{}6.738253e-02，2 次求值，err=\allowbreak{}7.39e-11 \checkmark \\
A & Bratu 打靶 Newton 的二次收敛 & |F\_\{k+1\}| $\le$ 0.2$\cdot$|F\_k|\textasciicircum{}2(\allowbreak{}牛顿法的标准二次界，只查 |F|$\ge$1e-4 的段) & 残差史 2.69e+01 -> 8.52e+00 -> 1.92e+00 -> 1.82e-01 -> 2.12e-03 -> 3.03e-07 -> 7.25e-13 \checkmark \\
B & 两条路线的灵敏度一致(相对) & $\le$ 1.00000e-09 & 0 \checkmark \\
A & 增长模态 vs 衰减模态(同一 S) & S=\allowbreak{}sinh(\allowbreak{}kL)/\allowbreak{}k 相同，但衰减模态要在指数小的分量上做相减 & S=\allowbreak{}5.8846e+15；增长模态 err=\allowbreak{}4.82e-10，衰减模态 err=\allowbreak{}1.73e+01，相差 3.59e+10 倍 \checkmark \\
A & 精确斜率 + 加密步长下的衰减模态误差 & 舍入被 e\textasciicircum{}\{kL\} 放大：误差 \textasciitilde{} eps$\cdot$e\textasciicircum{}\{kL\}，与 h 无关 & ns=\allowbreak{}250..8000(\allowbreak{}h: 8.0e-3 $\to$ 2.5e-4)误差 6.02e-01\textasciitilde{}2.60e+01，kL=\allowbreak{}40，eps$\cdot$exp(\allowbreak{}kL)=\allowbreak{}2.6e+01 \checkmark \\
A & Dirichlet 两端：整体阶 & 2 & 2.00078 \checkmark \\
A & Dirichlet 两端：内部阶 & 2 & 2.00078 \checkmark \\
B & 离散边界残差(Dirichlet 被钉死) & $\le$ 1.00000e-11 & $2.316\times 10^{-12}$ \checkmark \\
A & Dirichlet：整体阶 = 内部阶 = 2 & 二阶中心差商 + 精确 Dirichlet 行 & 整体 2.0008，内部 2.0008；误差 4.331e-03 -> 1.688e-05 \checkmark \\
A & 混合边界(边界格式 1 阶)：整体阶 & 1 & 0.99946 \checkmark \\
A & 混合边界(边界格式 2 阶)：整体阶 & 2 & 1.97556 \checkmark \\
A & 混合边界：整体阶与内部阶 & 整体阶 =\allowbreak{} min(\allowbreak{}内部阶,\allowbreak{} 边界格式阶) & 一阶边界：整体 0.9995 /\allowbreak{} 内部 0.9994；二阶边界：整体 1.9756 /\allowbreak{} 内部 1.9780 \checkmark \\
B & 离散边界残差(格式自己那一行) & $\le$ 1.00000e-12 & $7.886\times 10^{-14}$ \checkmark \\
A & 边界残差两个口径的差距(一阶边界格式) & 离散行残差 \textasciitilde{} eps；真实违反量 \textasciitilde{} h$\cdot$|u''| & n=\allowbreak{}64: 离散 3.3e-15 /\allowbreak{} 真实 1.9e-02 /\allowbreak{} 误差 1.9e-02；n=\allowbreak{}128: 离散 1.8e-14 /\allowbreak{} 真实 9.6e-03 /\allowbreak{} 误差 9.6e-03；n=\allowbreak{}256: 离散 7.9e-14 /\allowbreak{} 真实 4.8e-03 /\allowbreak{} 误差 4.8e-03；两者相差 6.1e+10 倍 \checkmark \\
B & 线性系统残差(LU 回代) & $\le$ 1.00000e-09 & $2.762\times 10^{-11}$ \checkmark \\
A & 二阶边界格式下真实边界残差的衰减阶 & 3 & 3.00035 \checkmark \\
A & 残差与误差的比(有限差分) & 残差 \textasciitilde{} eps$\cdot$||A||，误差由截断决定，两者量级无关 & 最大线性残差 2.76e-11，误差 4.45e-05(\allowbreak{}n=\allowbreak{}64)-> 7.73e-07(\allowbreak{}n=\allowbreak{}512)；最小比值 2.8e+04 \checkmark \\
B & 每行非零元个数(三对角) & $\le$ 3.05 & 2.9922 \checkmark \\
A & FD 矩阵的稀疏结构与消元代价 & 三对角：每行 3 个非零；稠密 LU \textasciitilde{} n\textasciicircum{}3/\allowbreak{}3，追赶法 \textasciitilde{} 5n & 每行非零 ['2.88',\allowbreak{} '2.97',\allowbreak{} '2.99']；n=\allowbreak{}512 时稠密/\allowbreak{}带状 =\allowbreak{} 17476.3(\allowbreak{}n=\allowbreak{}32: 68.3) \checkmark \\
B & 打靶 vs 闭式解 & $\le$ 1.00000e-09 & $1.529\times 10^{-12}$ \checkmark \\
B & 有限差分(n=320)vs 闭式解 & $\le$ 1.00000e-04 & $4.671\times 10^{-6}$ \checkmark \\
B & 谱配置(N=32)vs 闭式解 & $\le$ 1.00000e-11 & $1.110\times 10^{-15}$ \checkmark \\
A & Bratu：三条路线的误差与 Newton 迭代数 & u=\allowbreak{}4/\allowbreak{}(\allowbreak{}1+x)\textasciicircum{}2 & 打靶 1.53e-12(\allowbreak{}7 次 F 求值)；FD n=\allowbreak{}320 4.67e-06(\allowbreak{}5 步)；配置 N=\allowbreak{}32 1.11e-15(\allowbreak{}5 步)；互差界 4.67e-06 \checkmark \\
A & 非线性 FD 的 Newton 残差地板 & 线性系统残差地板 \textasciitilde{} eps$\cdot$||A|| \textasciitilde{} eps/\allowbreak{}h\textasciicircum{}2 & n=\allowbreak{}80: 残差 9.63e-12(\allowbreak{}eps/\allowbreak{}h\textasciicircum{}2=\allowbreak{}7.1e-13)；n=\allowbreak{}160: 残差 3.72e-11(\allowbreak{}eps/\allowbreak{}h\textasciicircum{}2=\allowbreak{}2.8e-12)；n=\allowbreak{}320: 残差 1.94e-10(\allowbreak{}eps/\allowbreak{}h\textasciicircum{}2=\allowbreak{}1.1e-11) \checkmark \\
C & 衰减模态打靶失效长度 L* & 1.83684 & 1.87488 \checkmark \\
C & 失效长度 L* 的两侧 & L* =\allowbreak{} ln(\allowbreak{}1/\allowbreak{}u)/\allowbreak{}k =\allowbreak{} 1.8368 & 二分 1.8749(\allowbreak{}区间 [1.8748,\allowbreak{} 1.8750]，13 次迭代，15 次谓词调用)；kL* =\allowbreak{} 37.50 \checkmark \\
C & 衰减模态打靶误差的指数律 & err \textasciitilde{} C$\cdot$e\textasciicircum{}\{kL\}(\allowbreak{}舍入被灵敏度放大)，半对数斜率 =\allowbreak{} 1 & kL=\allowbreak{}[5.0,\allowbreak{} 10.0,\allowbreak{} 15.0,\allowbreak{} 20.0,\allowbreak{} 25.0,\allowbreak{} 30.0] 对应 err=\allowbreak{}['1.19e-13',\allowbreak{} '1.92e-12',\allowbreak{} '3.54e-11',\allowbreak{} '8.21e-10',\allowbreak{} '2.34e-06',\allowbreak{} '5.24e-04']；拟合斜率 0.8927，R\textasciicircum{}2=\allowbreak{}0.95547 \checkmark \\
C & 同一格式两个窗口的实测阶 & 二阶格式的渐近阶 2；前渐近区会偏低 & n=\allowbreak{}16..128 得 1.7724；n=\allowbreak{}256..1024 得 1.9842 \checkmark \\
C & FD 矩阵条件数 \textasciitilde{} n\textasciicircum{}2 & cond\_inf \textasciitilde{} n\textasciicircum{}2(\allowbreak{}差分算子 \textasciitilde{} (\allowbreak{}1/\allowbreak{}h\textasciicircum{}2)$\cdot$三对角) & n=\allowbreak{}[32.0,\allowbreak{} 64.0,\allowbreak{} 128.0,\allowbreak{} 256.0] 对应 cond=\allowbreak{}['1.025e+03',\allowbreak{} '4.097e+03',\allowbreak{} '1.639e+04',\allowbreak{} '6.554e+04']；拟合指数 1.9996(\allowbreak{}R\textasciicircum{}2=\allowbreak{}1.00000) \checkmark \\
B & 变换-求值往返(值尺度 \textasciitilde{}8) & $\le$ 7.10543e-15 & $6.828\times 10^{-15}$ \checkmark \\
B & x\_\allowbreak{}j = cos(j pi/N)(精确到舍入) & $\le$ 1.00000e-15 & 0 \checkmark \\
A & Chebyshev 往返变换与节点定义 & 互逆；x\_j=\allowbreak{}cos(\allowbreak{}j pi/\allowbreak{}N)，j=\allowbreak{}0..N(\allowbreak{}降序、\allowbreak{}含端点) & 往返最大偏差 6.828e-15；节点最大偏差 0.000e+00(\allowbreak{}N=\allowbreak{}4..64) \checkmark \\
B & 一阶导数：系数空间 vs D\_\allowbreak{}N & $\le$ 1.00000e-10 & $1.883\times 10^{-13}$ \checkmark \\
B & 二阶导数：系数空间 vs D\_\allowbreak{}N\textasciicircum{}2 & $\le$ 1.00000e-08 & $5.877\times 10^{-11}$ \checkmark \\
A & 系数空间微分 vs 微分矩阵 & 系数空间求导 M1 与 D\_N 一致、\allowbreak{}M2 与 D\_N\textasciicircum{}2 一致(\allowbreak{}含 c\_0=\allowbreak{}c\_N=\allowbreak{}2 的因子) & N=\allowbreak{}24：一阶偏差 1.88e-13，二阶偏差 5.88e-11；再作用一次的自洽性偏差 2.55e+01 \checkmark \\
B & 对 m$\le$N 的 x\textasciicircum{}m 积分误差 & $\le$ 1.00000e-13 & $4.441\times 10^{-16}$ \checkmark \\
A & Clenshaw--Curtis 的代数精度 & $\int$\_\{-1\}\textasciicircum{}\{1\} x\textasciicircum{}m dx =\allowbreak{} 2/\allowbreak{}(\allowbreak{}m+1)(\allowbreak{}m 偶)，否则 0；代数精度 =\allowbreak{} N(\allowbreak{}N+2 处失效) & m$\le$N 的最大偏差 4.441e-16；N=\allowbreak{}4: m=\allowbreak{}N+2 处偏差 1.9e-02，sum(\allowbreak{}w)=\allowbreak{}2.00000000000000；N=\allowbreak{}8: m=\allowbreak{}N+2 处偏差 7.2e-05，sum(\allowbreak{}w)=\allowbreak{}2.00000000000000；N=\allowbreak{}16: m=\allowbreak{}N+2 处偏差 3.1e-08，sum(\allowbreak{}w)=\allowbreak{}2.00000000000000；N=\allowbreak{}32: m=\allowbreak{}N+2 处偏差 5.7e-14，sum(\allowbreak{}w)=\allowbreak{}2.00000000000000 \checkmark \\
B & ||D\_\allowbreak{}N||\_\allowbreak{}inf / N\textasciicircum{}2 - 1 & $\le$ 1.00000e-12 & $8.438\times 10^{-15}$ \checkmark \\
B & 行和(D\_\allowbreak{}N$\cdot$1 = 0) & $\le$ 1.00000e-10 & $2.274\times 10^{-13}$ \checkmark \\
B & D\_\allowbreak{}00 vs (2N\textasciicircum{}2+1)/6(相对) & $\le$ 1.00000e-13 & $1.082\times 10^{-14}$ \checkmark \\
B & ||D\_\allowbreak{}N\textasciicircum{}2||\_\allowbreak{}inf / (N\textasciicircum{}4/3) - 1 & $\le$ 0.02 & 0.015625 \checkmark \\
A & 微分矩阵的结构体检 & D 稠密(\allowbreak{}density=\allowbreak{}1)、\allowbreak{}D$\cdot$1=\allowbreak{}0、\allowbreak{}D\_00=\allowbreak{}(\allowbreak{}2N\textasciicircum{}2+1)/\allowbreak{}6、\allowbreak{}||D||\_inf=\allowbreak{}N\textasciicircum{}2、\allowbreak{}||D\textasciicircum{}2||\_inf$\approx$N\textasciicircum{}4/\allowbreak{}3 & N=\allowbreak{}8: ||D||=\allowbreak{}64(\allowbreak{}N\textasciicircum{}2=\allowbreak{}64)，行和 5.0e-16，||D\textasciicircum{}2||=\allowbreak{}1.344e+03；N=\allowbreak{}16: ||D||=\allowbreak{}256(\allowbreak{}N\textasciicircum{}2=\allowbreak{}256)，行和 4.4e-15，||D\textasciicircum{}2||=\allowbreak{}2.176e+04；N=\allowbreak{}32: ||D||=\allowbreak{}1024(\allowbreak{}N\textasciicircum{}2=\allowbreak{}1024)，行和 2.8e-14，||D\textasciicircum{}2||=\allowbreak{}3.492e+05；N=\allowbreak{}64: ||D||=\allowbreak{}4096(\allowbreak{}N\textasciicircum{}2=\allowbreak{}4096)，行和 2.3e-13，||D\textasciicircum{}2||=\allowbreak{}5.591e+06；最小密度 0.9877；D\_00 与 (\allowbreak{}2N\textasciicircum{}2+1)/\allowbreak{}6 最大相对偏差 1.08e-14 \checkmark \\
B & 配置法内点残差 & $\le$ 1.00000e-08 & $5.157\times 10^{-12}$ \checkmark \\
A & 配置法的残差三口径 & 配置点内：ODE 残差 =\allowbreak{} 舍入；节点之间：残差 \textasciitilde{} 误差$\cdot$N\textasciicircum{}2(\allowbreak{}插值误差的导数) & N=\allowbreak{}8: 内点残差 1.7e-13，细网格残差 3.8e-01，误差 1.8e-04；N=\allowbreak{}12: 内点残差 1.9e-12，细网格残差 6.1e-04，误差 6.0e-08；N=\allowbreak{}16: 内点残差 2.3e-12，细网格残差 1.0e-07，误差 3.5e-12；N=\allowbreak{}20: 内点残差 5.2e-12，细网格残差 2.9e-11，误差 8.5e-15；细网格残差/\allowbreak{}误差 最小 2104 倍 \checkmark \\
A & 精确解的谱微分残差地板 & 微分放大舍入 \textasciitilde{}N\textasciicircum{}4$\cdot$u(\allowbreak{}谱微分矩阵 ||D\textasciicircum{}2||\textasciitilde{}N\textasciicircum{}4) & N=\allowbreak{}20: 残差 6.82e-12(\allowbreak{}N\textasciicircum{}4u=\allowbreak{}1.8e-11)，插值误差 3.33e-15；N=\allowbreak{}24: 残差 2.10e-11(\allowbreak{}N\textasciicircum{}4u=\allowbreak{}3.7e-11)，插值误差 2.44e-15；N=\allowbreak{}32: 残差 1.08e-10(\allowbreak{}N\textasciicircum{}4u=\allowbreak{}1.2e-10)，插值误差 2.66e-15；实测/\allowbreak{}预测 最大比例 0.93 \checkmark \\
B & 配置法与 tau 法的解之差 & $\le$ 1.00000e-10 & $2.677\times 10^{-11}$ \checkmark \\
A & 配置法 vs tau 法 & 同一多项式空间、\allowbreak{}同一方程：两条路线应当给同一个解 & N=\allowbreak{}16: 配置 3.5e-12 /\allowbreak{} tau 2.7e-11 /\allowbreak{} 互差 2.7e-11；N=\allowbreak{}20: 配置 8.5e-15 /\allowbreak{} tau 6.9e-15 /\allowbreak{} 互差 9.7e-15；N=\allowbreak{}24: 配置 7.5e-15 /\allowbreak{} tau 5.5e-15 /\allowbreak{} 互差 9.1e-15；N=\allowbreak{}32: 配置 1.9e-14 /\allowbreak{} tau 5.4e-15 /\allowbreak{} 互差 2.1e-14 \checkmark \\
B & 谱配置的边界残差 & $\le$ 1.00000e-11 & $1.990\times 10^{-13}$ \checkmark \\
A & 谱配置的边界残差与误差同量级 & 边界行是 D\_N 的一行，解出来时行残差 \textasciitilde{} 舍入(\allowbreak{}与内点同量级) & smooth N=\allowbreak{}16: 离散=\allowbreak{}5.9e-16 真实=\allowbreak{}5.9e-16 误差=\allowbreak{}3.5e-12；smooth N=\allowbreak{}24: 离散=\allowbreak{}4.6e-15 真实=\allowbreak{}4.6e-15 误差=\allowbreak{}7.5e-15；smooth N=\allowbreak{}32: 离散=\allowbreak{}8.8e-15 真实=\allowbreak{}8.8e-15 误差=\allowbreak{}1.9e-14；neumann N=\allowbreak{}16: 离散=\allowbreak{}8.2e-14 真实=\allowbreak{}8.2e-14 误差=\allowbreak{}1.1e-13；neumann N=\allowbreak{}24: 离散=\allowbreak{}8.5e-14 真实=\allowbreak{}8.5e-14 误差=\allowbreak{}6.2e-15；neumann N=\allowbreak{}32: 离散=\allowbreak{}2.0e-13 真实=\allowbreak{}2.0e-13 误差=\allowbreak{}4.2e-13 \checkmark \\
B & 聚集比 / N - 2/pi & $\le$ 0.001 & $3.196\times 10^{-5}$ \checkmark \\
A & 节点聚集的几何量 & 间距比 \textasciitilde{} (\allowbreak{}2/\allowbreak{}pi)N；密度 rho(\allowbreak{}x)=\allowbreak{}N/\allowbreak{}(\allowbreak{}pi sqrt(\allowbreak{}1-x\textasciicircum{}2))，rho(\allowbreak{}0.9)/\allowbreak{}rho(\allowbreak{}0)=\allowbreak{}1/\allowbreak{}sqrt(\allowbreak{}0.19)=\allowbreak{}2.294 & N=\allowbreak{}8: 最大/\allowbreak{}最小间距 =\allowbreak{} 5.027(\allowbreak{}/\allowbreak{}N =\allowbreak{} 0.6284)；N=\allowbreak{}16: 最大/\allowbreak{}最小间距 =\allowbreak{} 10.153(\allowbreak{}/\allowbreak{}N =\allowbreak{} 0.6346)；N=\allowbreak{}32: 最大/\allowbreak{}最小间距 =\allowbreak{} 20.355(\allowbreak{}/\allowbreak{}N =\allowbreak{} 0.6361)；N=\allowbreak{}64: 最大/\allowbreak{}最小间距 =\allowbreak{} 40.735(\allowbreak{}/\allowbreak{}N =\allowbreak{} 0.6365)；N=\allowbreak{}128: 最大/\allowbreak{}最小间距 =\allowbreak{} 81.483(\allowbreak{}/\allowbreak{}N =\allowbreak{} 0.6366)；N=\allowbreak{}64 时 rho(\allowbreak{}0.9)/\allowbreak{}rho(\allowbreak{}0)=\allowbreak{}2.294157(\allowbreak{}解析 2.294157) \checkmark \\
B & Bratu(x$\in$[0,1])的谱配置误差 & $\le$ 1.00000e-11 & $2.220\times 10^{-15}$ \checkmark \\
A & 区间映射的伸缩因子 s=2/(b-a) & d/\allowbreak{}dx =\allowbreak{} s$\cdot$d/\allowbreak{}dxi，d\textasciicircum{}2/\allowbreak{}dx\textasciicircum{}2 =\allowbreak{} s\textasciicircum{}2$\cdot$d\textasciicircum{}2/\allowbreak{}dxi\textasciicircum{}2；[0,\allowbreak{}1] 上 s=\allowbreak{}2 & 正确伸缩：误差 2.22e-15；漏掉伸缩：偏差 7.59e-01(\allowbreak{}s=\allowbreak{}2 时算子整体差 4 倍) \checkmark \\
C & 指数收敛 vs 幂律收敛的判别 & 解析解 err\textasciitilde{}e\textasciicircum{}\{-lambda N\}(\allowbreak{}比值递增)；C\textasciicircum{}s 解 err\textasciitilde{}N\textasciicircum{}\{-s\}(\allowbreak{}比值 $\to$ 常数) & smooth(\allowbreak{}N=\allowbreak{}8..18)：斜率 -2.2568，R\textasciicircum{}2=\allowbreak{}0.99636，误差比 ['39.3',\allowbreak{} '76.0',\allowbreak{} '196.7',\allowbreak{} '88.0']；abs3(\allowbreak{}N=\allowbreak{}8..32)：p=\allowbreak{}1.9729，R\textasciicircum{}2=\allowbreak{}0.999996，误差比 ['1.56',\allowbreak{} '1.43',\allowbreak{} '1.36',\allowbreak{} '1.30',\allowbreak{} '1.26',\allowbreak{} '1.23',\allowbreak{} '1.21',\allowbreak{} '1.19',\allowbreak{} '1.17',\allowbreak{} '1.16',\allowbreak{} '1.15',\allowbreak{} '1.14'] \checkmark \\
C & C\textasciicircum{}s 解的配置率 = 插值率 - 1 & 微分吃掉一阶光滑度：配置率 =\allowbreak{} 插值率 - 1(\allowbreak{}实测 |x|\textasciicircum{}3: 3.00$\to$1.99；|x|\textasciicircum{}5: 5.01$\to$4.02) & abs3: 配置解 p=\allowbreak{}1.987，插值 p=\allowbreak{}3.002，差 1.015(\allowbreak{}理论 1)；abs5: 配置解 p=\allowbreak{}4.023，插值 p=\allowbreak{}5.012，差 0.989(\allowbreak{}理论 1) \checkmark \\
C & 解析算例的舍入地板 & 地板 \textasciitilde{} N\textasciicircum{}4$\cdot$u(\allowbreak{}微分放大舍入)，N$\approx$18 起到达 & 地板 7.550e-15(\allowbreak{}N$\ge$18)，最小误差 24 处 7.55e-15；N=\allowbreak{}64 时 9.26e-14(\allowbreak{}比 N=\allowbreak{}24 的 7.55e-15 更大) \checkmark \\
C & 微分矩阵的条件数增长 & kappa+ \textasciitilde{} N\textasciicircum{}2(\allowbreak{}sigma\_min+ 有界)，朴素 kappa 由零空间(\allowbreak{}常数)决定，无意义 & kappa+ =\allowbreak{} ['35.1',\allowbreak{} '74.8',\allowbreak{} '130.0',\allowbreak{} '286.2',\allowbreak{} '503.9',\allowbreak{} '1123.2',\allowbreak{} '1987.9']；拟合指数 1.946(\allowbreak{}R\textasciicircum{}2=\allowbreak{}0.99989)；朴素 kappa 最大 5.93e+18，比 kappa+ 大 3.0e+15 倍 \checkmark \\
C & 陡峭解析解所需的分辨率 & tanh(\allowbreak{}kx) 在带形内解析 =\allowbreak{}> 仍指数收敛，但速率 \textasciitilde{} e\textasciicircum{}\{-c$\cdot$N/\allowbreak{}k\} & err<1e-10 的最小 N：k=\allowbreak{}3 -> 48，k=\allowbreak{}6 -> 96，k=\allowbreak{}10 -> None(\allowbreak{}窗口上限 128)；半对数斜率 k=\allowbreak{}3: -0.589，k=\allowbreak{}6: -0.276，k=\allowbreak{}10: -0.111 \checkmark \\
\end{longtable}
